\problem{Basler Problem}
    Consider the famous series
    \begin{equation}
        \sum_{k=1}^\infty k^{-2} = \frac{\pi^2}{6}
    \end{equation}
    known as the Basler problem. We would like to approximate 
    this series by finitely many summands.
    \begin{enumerate}
        \item For a given finite expression $1 + \frac{1}{2^2} + \dots + \frac{1}{n^2}$ will it make a difference if we sum the number from small to large or from large to small on the computer at finite precision?
        \item Implement both the summation from large to small and small to large in single precision (use np.float32) and compare the results to the ground truth, $\frac{\pi^2}{6}$.
    \end{enumerate}
\begin{solutionblock}
    \part
    If we sum the terms just as the formula suggests 
    from large to small, at some point, the small changes 
    will not be resolved anymore - so better
    sum up from small to large.
    \part
    See Code-Snippet \ref{code:basel}.

    \begin{codebox}
        \begin{minted}{python3}
        import numpy as np
        import math

        def basel_sum_small_to_large_f32(n):
            s = np.float32(0.0)
            for k in range(n, 0, -1):
                s += np.float32(1.0) / np.float32(k * k)
            return s

        def basel_sum_large_to_small_f32(n):
            s = np.float32(0.0)
            for k in range(1, n + 1):
                s += np.float32(1.0) / np.float32(k * k)
            return s

        # Ground truth (double precision)
        exact = math.pi**2 / 6

        for n in [10**3, 10**4, 10**5, 10**6]:
            s_sl = basel_sum_small_to_large_f32(n)
            s_ls = basel_sum_large_to_small_f32(n)

            err_sl = abs(float(s_sl) - exact)
            err_ls = abs(float(s_ls) - exact)

            print(f"n = {n:>7}")
            print(f"  small -> large (float32): {s_sl:.7f}, error = {err_sl:.3e}")
            print(f"  large -> small (float32): {s_ls:.7f}, error = {err_ls:.3e}")
            print()
        \end{minted}
        \caption{Python code for approximating $\sum_{k=1}^\infty k^{-2}$ using $n$ summands. Summation from
        large to small and from small to large in finite precision are compared.}
        \label{code:basel}
    \end{codebox}
    
\end{solutionblock}